\begin{abstract}
Can intermediate supervision teach a recurring transition quickly when final-answer supervision supplies exponentially weak credit? We study modular counting in a concrete two-block causal softmax-attention/ReLU Transformer. Near a small initialization, the outcome loss cannot fall by more than $O(\varepsilon^2\lambda^n)$, and a general escape-time lemma turns this loss floor into an $\Omega(\lambda^{-n})$ plateau of gradient flow on every parameter. In an aligned supplied-routing family, process supervision develops persistent positive margins on every recurring transition in $O(n^{5/2}\log n)$ time; from a common seed this gives a matched learning-time separation. We add a conditional theorem for the full Transformer with all embedding, attention, MLP and readout parameters trained: if the loss gradient keeps an inverse-polynomial alignment with a bounded direction, the process loss falls in polynomial time, and a first-error argument turns this into correct generation of the entire autoregressive trace with high probability. A readout-separability condition on the evolving hidden features implies the alignment. The plateau holds from a Gaussian seed with a weak perturbation, zero readout and sign MLP biases; proving the process alignment from that seed is open, and we state the exact matrix dynamics, a conditional cancellation lemma and the measurements that test it. Counting and pretrained Qwen2.5-0.5B LoRA experiments support the qualitative mechanism; a five-object BIG-Bench Hard reversal limits its empirical scope.
\end{abstract}

\section{Introduction}
A final answer can conceal many repeated state updates. A running trace exposes those updates as local prediction problems. The question is whether this changes optimization time within the same Transformer, rather than merely changing what the architecture can represent. We study counting modulo $K$: each independent letter either keeps the state or advances it by one. The answer space stays fixed as the computation grows.

The paper has four levels, each with a different status.
\begin{enumerate}[label=(\arabic*),leftmargin=*,itemsep=1pt,topsep=2pt]
\item \emph{Representation.} Explicit two-block executors compute both the direct final count and the full running trace at linear width (\cref{prop:represent}). This is existence, not learnability.
\item \emph{Outcome barrier} (\cref{thm:outcome}, proved). Near a small initialization the outcome loss has a floor $\log K-\Delta_n$ with $\Delta_n=O(\varepsilon^2\lambda^n)$. Gradient flow on every parameter then needs time $\Omega(\lambda^{-n})$ to leave the neighbourhood.
\item \emph{Matched separation} (\cref{thm:step,cor:matched}, proved). In an aligned family with supplied routing and tied query--key scalars, process supervision reaches persistent positive margins in $O(n^{5/2}\log n)$ time, while outcome supervision from the same seed and in the same family is still on its plateau:
\begin{equation}\label{eq:headline}
\boxed{T_{\mathrm{step}}=O(n^{5/2}\log n),\qquad T_{\mathrm{out}}=\Omega(\lambda^{-n}).}
\end{equation}
\item \emph{Full training} (\cref{sec:full}). With every tensor trained, persistent gradient alignment implies polynomial-time learning and correct free-running generation (\cref{thm:process}, conditional). The outcome plateau holds from a Gaussian seed with a weak perturbation (\cref{prop:random-plateau}, proved); the alignment from that seed is open.
\end{enumerate}
$T_{\mathrm{out}}$ is the duration of an initial plateau, measured by displacement, cross-entropy and expected probability of the correct answer. It is not a lower bound for reaching arbitrary accuracy, and it says nothing about permanent failure. None of the learning results identifies which heads or MLP units carry the computation; they concern predictions.

The two mechanisms are different in kind. The outcome argument splits into a model-specific bound on answer credit, which Fourier contraction makes exponentially small, and an architecture-free lemma: squared speed along gradient flow equals the rate of loss decrease, so a small loss floor forces slow escape. The process argument uses the displayed state, which leaves a rule signal of order one at every transition. Small-initialization counting runs measure displacement times; other counting runs measure rule learning; pretrained Qwen LoRA experiments test qualitative transfer, including a case where outcome supervision performs better. None of these results proves either bound.

\section{Modular counting and Transformer setup}\label{sec:setup}\label{sec:setting}
\begin{definition}[Task and contraction]\label{def:task}\label{def:lambda}
Fix $K\ge2$, a finite alphabet $\Sigma$, and a query letter $w^\star$ with probability $p\in(0,1)$. Independently sample $s_0$ uniformly from $\mathbb Z_K$ and letters $w_1,\ldots,w_n$. Put $b_t=\mathbf1[w_t=w^\star]$ and
\begin{equation}
s_t=s_{t-1}+b_t\pmod K,\qquad y=s_n,\qquad
\lambda=\max_{1\le j<K}|1-p+p e^{2\pi\mathrm i j/K}|.
\label{eq:lambda}
\end{equation}
For $p\ne1/2$, $0<\lambda<1$. In the binary case $\lambda=|1-2p|$.
\end{definition}
The prompt displays $s_0$, the query, and all letters. Outcome format predicts only the final count. Process format appends the running states $s_1,\ldots,s_n$, a colon, the final count and EOS: a continuation $y_1,\ldots,y_M$ of $M=n+3$ tokens. Training uses gold previous tokens. At the position $i_t=n+1+t$ predicting $s_t$, $t\ge2$, the displayed token is $s_{t-1}$; the current letter is at $j_t=1+t$.

\begin{definition}[Population objectives]\label{def:losses}
With $z_t(\Theta)\in\mathbb R^K$ the count logits predicting $s_t$, $z_{\mathrm{ans}}(\Theta)$ the outcome-answer logits, and $p_\Theta(\cdot\mid P,y_{<m})$ the next-token distribution of the process format,
\begin{equation}\label{eq:objectives}
\begin{gathered}
\Lout=\E[-\log\softmax(z_{\mathrm{ans}})_y],\qquad
\Lstep=\frac1{n-1}\sum_{t=2}^n\E[-\log\softmax(z_t)_{s_t}],\\
\Lproc=\frac1M\sum_{m=1}^M\E\bigl[-\log p_\Theta(y_m\mid P,y_{<m})\bigr].
\end{gathered}
\end{equation}
$\Lstep$ covers the recurring transitions only and is the objective of \cref{thm:step}. $\Lproc$ covers the whole continuation, including the first state and the formatting tokens, and is the objective of \cref{thm:process}. They are different objectives.
\end{definition}

\paragraph{Overall architecture.}\label{ass:arch}
Each of two residual blocks has three causal softmax-attention heads and a ReLU MLP. There is no LayerNorm. A head may use a fixed input-independent causal submask. With $h_i^0=E(\mathrm{token}_i)+P_i$, each block computes
\begin{align}
\bar h_i^\ell&=h_i^{\ell-1}+\sum_{a=1}^3\sum_{j\in\mathcal M_a(i)}\omega_{ij}^{\ell,a} W_V^{\ell,a}h_j^{\ell-1},\nonumber\\
\omega_{ij}^{\ell,a}&=\softmax_{j\in\mathcal M_a(i)}\left(\frac{(W_Q^{\ell,a}h_i^{\ell-1})^\top W_K^{\ell,a}h_j^{\ell-1}}{\sqrt{d_{\mathrm h}}}+\beta_{i-j}^{\ell,a}\right),\label{eq:block}\\
h_i^\ell&=\bar h_i^\ell+W_{\mathrm{out}}^\ell\relu(W_{\mathrm{in}}^\ell\bar h_i^\ell+b^\ell),\qquad z_i=Vh_i^2.\nonumber
\end{align}
The value map includes the map back to the residual stream. The executor construction in \cref{app:architectures} specifies every register and nonzero tensor. At fixed $K,|\Sigma|$, head dimension, residual width and MLP width are of order $n$.

\paragraph{Notation.} $\Theta$ denotes all parameters of \cref{eq:block}: embeddings, the $18$ query, key and value maps, both MLPs and their biases, the readout $V$ and the relative biases $\beta$. Unless stated otherwise, training is Euclidean population gradient flow $\dot\Theta=-\nabla_\Theta\mathcal L$ on all of them. $\varepsilon$ is the initialization scale, and the \emph{outcome credit} is $\outcredit(\Theta)=\E\langle z_{\mathrm{ans}},e_y-u\rangle$ with $u$ uniform on $\mathbb Z_K$. Query and key matrices are always written $W_Q,W_K$; the letter $K$ is reserved for the modulus.

\section{Two explicit algorithms in one Transformer}\label{sec:results}\label{sec:executors}
Before asking what training finds, we fix what the architecture can do. Both computations below use the same two blocks and the same two primitives: a head that routes to (or averages over) intended positions, and a ReLU lookup that applies a modular update.
\begin{center}\small
\fitwidth{\begin{tabular}{@{}l@{\hspace{1.2em}}l@{\hspace{1.2em}}l@{}}
\toprule
& \textbf{process executor} & \textbf{outcome executor}\\
\midrule
block 1 (shared) & \multicolumn{2}{l}{retrieve the query $w^\star$ at every letter; MLP computes $b_t=\mathbf 1[w_t=w^\star]$}\\
block 2, attention & retrieve $s_{t-1}$ and $b_t$ (two tokens) & average $b_1,\ldots,b_n$, retrieve $s_0$ ($n+1$ tokens)\\
block 2, MLP & $s_t=s_{t-1}+b_t\bmod K$ & $s_n=s_0+\sum_tb_t\bmod K$\\
output & the whole trace $s_1,\ldots,s_n,{:},s_n,\mathrm{EOS}$ & the final count\\
\bottomrule
\end{tabular}}
\end{center}

\begin{proposition}[Exact Outcome and Step Executors Exist in the Same Two-Block Transformer]\label{prop:represent}
For every $n\ge2$, $K\ge2$ and finite $\Sigma$, the architecture of \cref{eq:block} with full causal masks, three heads per block, $d_{\mathrm h}=2n+6$, $d_{\mathrm{res}}=2|\Sigma|+4K+d_{\mathrm h}+3$, and MLP width $\max\{|\Sigma|,3K(n+1),2n+1\}$ admits separate parameter settings that greedily generate the entire running-count continuation and the direct final-count continuation on every prompt. Every required prediction has strictly positive logit margin.
\end{proposition}
The proof is in \cref{app:architectures}; exhaustive finite verification is in \cref{app:verify}.
The construction (\cref{app:architectures}) is modular. A routing lemma bounds the attention mass that a finite score gap $\Omega_n=\max\{16,3\log(8(n+3))\}$ leaks outside the intended sources by $\omega<1/[256(n+3)^2]$; a lookup lemma shows that the units $\relu(x[a]+\mathbf 1[b=c]-\frac32)$ with output weights $2e_{a+c}$ compute $e_{s+b}$ exactly and stay exact under small input errors; a tent-function lookup does the same for $s_0+m$. Propagating the leak gives a uniform margin: at every valid prompt and every supervised position, including the first state and the formatting tokens, $z_{\mathrm{correct}}-\max_{k\ne\mathrm{correct}}z_k>19$ (\cref{prop:margins}). Both programs therefore have the same architectural access. What differs is the credit their objectives provide: the process program's block-2 inputs are two tokens, the outcome program's are an aggregate over all $n$ letters. The next two sections show what this does to training.


\section{Outcome credit and the exponential plateau}\label{sec:credit}\label{sec:outcome}
The final count depends on all $n$ independent increments. Near small initialization, a Transformer reaches this dependence only through high-order terms of its attention-score expansion, and a term that reads fewer than $n$ increments receives exponentially attenuated credit from the final-answer label. This weak credit has a direct optimization consequence: if the model cannot lower the outcome loss appreciably near its initialization, gradient flow cannot move far there, because its squared speed is exactly the rate of loss decrease.

\begin{theorem}[Outcome credit barrier]\label{thm:outcome}\label{thm:outcome-end}
Fix $K\ge2$ and $p\in(0,1)\setminus\{1/2\}$. Let the seed $\Theta_0=\varepsilon\tilde\Theta$ satisfy the small-initialization conditions of \cref{ass:init} (scaled operator norms at most $a_0$, value embedding nondegenerate, MLP biases bounded away from zero), with zero readout and zero relative biases. There are $c_{\mathrm{init}},c_\lambda,C_2>0$, depending only on $K$, $a_0$, $c_{\mathrm{emb}}$ and the block/head counts, such that if
\begin{equation}\label{eq:small}
0<\varepsilon\le\min\{1/2,\ c_{\mathrm{init}}\min(1,N_b^{-1}),\ c_\lambda\lambda\},
\end{equation}
then, with $r=r_0\varepsilon$ and $\Delta_n=C_2\varepsilon^2\lambda^n$:
\begin{thmparts}
\item \emph{Loss floor.} $\Lout(\Theta)\ge\log K-\Delta_n$ whenever the parameters other than the relative biases are within $2r$ of the seed, for arbitrary relative biases.
\item \emph{Plateau.} Along gradient flow of $\Lout$ on every parameter, for $0\le t\le T_{\mathrm{out}}:=r^2/(2\Delta_n)=r_0^2\lambda^{-n}/(2C_2)$,
\begin{equation}\label{eq:plateau}
\begin{gathered}
\|\Theta(t)-\Theta_0\|\le r,\qquad
\log K-\Delta_n\le\Lout(\Theta(t))\le\log K,\\
\bigl|\E[\softmax(z_{\mathrm{ans}})_y]-\tfrac1K\bigr|\le\sqrt{\Delta_n/2}.
\end{gathered}
\end{equation}
\end{thmparts}
\end{theorem}

\paragraph{One identity.} For one prompt, $\KL(u\|p)=\log\sum_ae^{z_a}-\log K-\langle u,z\rangle$. Averaging,
\begin{equation}\label{eq:main-identity}
\Lout(\Theta)=\log K+\E\bigl[\KL(u\,\|\,p_\Theta)\bigr]-\outcredit(\Theta)\ \ge\ \log K-\outcredit(\Theta).
\end{equation}
So the outcome loss can fall below $\log K$ only as far as the credit allows.

\paragraph{Small credit.} Fix the query and $s_0$. For $S\subseteq[n]$, Fourier conditioning gives
\begin{equation}\label{eq:credit-degree}
\|\E[H^\top(e_y-u)]\|\le\lambda^{n-|S|}\|H\|_{L^2}
\end{equation}
for every square-integrable matrix $H$ measurable in $w_S$: each omitted increment multiplies the $j$th Fourier mode by $1-p+pe^{2\pi\mathrm ij/K}$ (\cref{lem:target-contraction}). Inside a fixed-ReLU neighbourhood, scale every content score by $\zeta$ and expand at $\zeta=0$. A term of score order $k$ reads at most $2k+1$ letters and has size $C\varepsilon^2(C_{\mathrm{att}}\varepsilon^4)^k$; the zero-order network is affine in individual inputs. These are estimates for \cref{eq:block} itself, from the softmax derivative recursion and two-block composition. Summing \cref{eq:credit-degree} over the expansion gives $|\outcredit(\Theta)|\le\Delta_n=C_2\varepsilon^2\lambda^n$ on the neighbourhood (\cref{app:outcome}), which with \cref{eq:main-identity} is \cref{thm:outcome}(i).

\paragraph{Slow escape.} The remaining step does not use the architecture.
\begin{lemma}[A local loss floor implies a plateau]\label{lem:escape}
Let $F$ be continuously differentiable near the closed ball $\bar B(\Theta_0,r)$, let $\dot\Theta=-\nabla F(\Theta)$ with $\Theta(0)=\Theta_0$, and suppose $F\ge F(\Theta_0)-\Delta$ on $\bar B(\Theta_0,r)$. Then the exit time satisfies $T_{\mathrm{exit}}\ge r^2/\Delta$, and $\|\Theta(t)-\Theta_0\|^2\le t\Delta$ for $t\le T_{\mathrm{exit}}$.
\end{lemma}
\begin{proof}
Before exit, $\frac{d}{dt}F(\Theta(t))=-\|\dot\Theta\|^2$, so by Cauchy--Schwarz
$\|\Theta(t)-\Theta_0\|^2\le t\int_0^t\|\dot\Theta\|^2=t\,[F(\Theta_0)-F(\Theta(t))]\le t\Delta$. At exit the left side is $r^2$.
\end{proof}
With $F=\Lout$, $F(\Theta_0)=\log K$ (zero readout), radius $r=r_0\varepsilon$ and $\Delta=\Delta_n$, the lemma gives $T_{\mathrm{exit}}\ge r_0^2\lambda^{-n}/C_2$, and \cref{eq:main-identity} with Pinsker's inequality bounds the assigned probability. For the restricted flow, apply the lemma in $\alpha$ coordinates to $F\circ\iota$ on the ball of radius $r/\sqrt K$, which costs the factor $K$ in \cref{cor:matched}. In short:
\[
\text{Fourier contraction}\ \Longrightarrow\ \text{small credit}\ \Longrightarrow\ \text{local loss floor}\ \Longrightarrow\ \text{slow escape}.
\]

\section{Polynomial process learning}\label{sec:process}
\subsection{A displayed state leaves a local signal}
At a transition the displayed count $s_{t-1}$ already predicts $s_t$ with probabilities $(1-p,p)$. What it does not predict, the \emph{innovation} $\chi^{js_t}-\rho_j\chi^{js_{t-1}}$ in Fourier mode $j$ (with $\chi=e^{2\pi\mathrm i/K}$ and $\rho_j=1-p+p\chi^j$), has variance $1-|\rho_j|^2\ge1-\lambda^2$ (\cref{eq:move-variance}), independent of $n$. A feature that reads two tokens, the displayed count and the current letter, therefore carries order-one credit: the process loss pays the contraction once, not $n$ times.

\subsection{The aligned family}
\paragraph{The structured seed and the restricted family.}\label{sec:seed}
\Cref{thm:step} uses a seed in which only head 1 of block 2 is active. Token embeddings of scale $\varepsilon$ write the count and the match bit into separate directions; its mask at $i_t$ admits $i_t$ and $j_t$; its maps are $W_Q=q\Pi_Q$, $W_K=k\Pi_K$, $W_V=v\Pi_V$ for fixed matrices $\Pi_\bullet$; all other heads and MLP weights vanish and every MLP bias is $\pm\varepsilon$ (\cref{app:joint}). The \emph{restricted coordinates} are a count readout $U$, a gate readout $W=w(I-R)$ with $Re_s=e_{s+1}$, and the two scalars $q,k$:
\begin{equation}\label{eq:metric}
\alpha=(U,W,q,k),\qquad \Theta=\iota(\alpha),\qquad \|\alpha\|^2=\|U\|_F^2+\|W\|_F^2+q^2+k^2.
\end{equation}
The restricted flow is $\dot\alpha=-\nabla_\alpha\mathcal L(\iota(\alpha))$; everything else is fixed. It is \emph{not} the unrestricted flow on $\Theta$, and $\iota$ is not an isometry: $\|\delta\alpha\|^2\le\|\iota\,\delta\alpha\|^2\le K\|\delta\alpha\|^2$ (\cref{app:match}). The initialization is
\begin{equation}\label{eq:init}
U_0=W_0=0,\quad \varepsilon/2\le|q_0|,|k_0|,|v|\le2\varepsilon,
\quad q_0k_0(1-2p)<0.
\end{equation}
The last condition is \emph{alignment}: the initial gate points the useful way. A random choice of the two signs is aligned with probability $\frac12$.

\begin{theorem}[Process learning in the aligned family]\label{thm:step}
Under the conditions of \cref{thm:outcome} for the structured seed, with $h^\star=|\log((1-p)/p)|$ and $p_{\min}=\min(p,1-p)$, let $\alpha(t)$ follow the restricted flow of $\Lstep$ from \cref{eq:init}. There is $t_3\le T_{\mathrm{step}}$ after which every recurring transition $t\ge2$ has positive margin against all $K-1$ wrong counts, persistently, and $\Lstep\to0$. The explicit bound \cref{eq:focused-time} (\cref{app:match}) satisfies
\[
T_{\mathrm{step}}=O\!\left(\frac{Kh^\star}{p_{\min}(1-2p_{\min})\varepsilon^2}+\frac{\sqrt{d_{\mathrm h}}}{p_{\min}\min\{1,h^\star\}\varepsilon^4}\log\frac{d_{\mathrm h}}{\varepsilon}+\frac{K\log K}{p_{\min}^2\varepsilon^4}\right),
\]
which is $O(n^{5/2}\log n)$ at executor width $d_{\mathrm h}\asymp n$ and $\varepsilon\asymp n^{-1/2}$.
\end{theorem}

\begin{corollary}[Matched separation]\label{cor:matched}\label{thm:matched}
Run the restricted flows of $\Lout$ and $\Lstep$ from the same structured seed \cref{eq:init}, in the same family and metric \cref{eq:metric}. The outcome conclusions \cref{eq:plateau} hold for $t\le T_{\mathrm{out}}^\alpha:=r_0^2\lambda^{-n}/(2KC_2)$, the factor $K$ coming from $\iota$. If $T_{\mathrm{step}}<T_{\mathrm{out}}^\alpha$, both \cref{thm:step} and \cref{eq:plateau} hold on $[T_{\mathrm{step}},T_{\mathrm{out}}^\alpha]$. With $\varepsilon=c_*n^{-1/2}$, widths of order $n$ and fixed $K,p$, this holds for all sufficiently large $n$, which is \cref{eq:headline}.
\end{corollary}
The two flows share the seed and the trainable directions, not trajectories or teacher-forced inputs. At a finite horizon the exact inequality $T_{\mathrm{step}}<T_{\mathrm{out}}^\alpha$ is required; quoting orders does not establish it. The positive theorem fixes values, supplies routing and ties query--key scalars; the barrier does not need these restrictions.

In the aligned family this is turned into a rate. The logit reduces exactly to a count readout, a gate readout and the gate strength $a=\varepsilon^2qk/\sqrt{d_{\mathrm h}}$, and their combination that uses the letter, the \emph{rule contrast}, does not cancel with the count prior (\cref{app:rule-contrast}). For $p<\frac12$ the state prior moves the count readout first; alignment makes the same signal grow the gate readout the right way; a balance law keeps it aligned, so the gate strength grows exponentially to order one; readout descent then gives positive margins. Every competitor logit is controlled (\cref{app:joint}).

\paragraph{What the observables mean.}
Positive transition margins imply correct greedy transitions given the gold previous state. The outcome bound controls the average probability assigned to the correct answer; it does not bound greedy accuracy, which can be sensitive to arbitrarily small logit differences.

\subsection{Training every parameter}\label{sec:full}
We now drop all restrictions: every tensor of \cref{eq:block} trains by the same Euclidean flow, nothing is frozen, tied, projected or rescaled, and masks are full causal masks.

\paragraph{A random seed.} Let
\begin{equation}\label{eq:rand-seed}
\Theta_0=\varepsilon G+\delta A .
\end{equation}
In $G$, every weight and embedding block other than the readout has i.i.d.\ $\mathcal N(0,1/m)$ entries, with $m$ its larger dimension; the readout and relative biases of $G$ are zero, and its scaled MLP biases are fixed signs. The perturbation $A$ is deterministic and task-aligned (for example weak positional or match-sensitive directions), vanishes on the readout and biases, and satisfies $\delta\|A_{\mathrm B}\|_{\mathrm{op}}\le\varepsilon/4$ on every block $\mathrm B$. Since the readout is zero, every initial prediction is uniform: the seed does not count.

\begin{lemma}[The random seed is a small initialization]\label{lem:seed}
Let $N$ be the number of parameter blocks, $m_{\min}$ the smallest larger-dimension among them, and suppose the residual width is at least $16$ times the vocabulary size. With probability at least $1-2(N+1)e^{-m_{\min}/32}$ over $G$, the seed \cref{eq:rand-seed} satisfies \cref{ass:init} with $a_0=4$, $c_{\mathrm{emb}}=1/4$, $r_0=1/16$ and $N_b=\sqrt w$.
\end{lemma}

\begin{proposition}[Outcome plateau from the random seed]\label{prop:random-plateau}
On the event of \cref{lem:seed}, with $\varepsilon$ as in \cref{eq:small} and $N_b=\sqrt w$, the conclusions of \cref{thm:outcome} hold for the flow of $\Lout$ on every parameter from \cref{eq:rand-seed}, with $r_0=1/16$. At executor widths the event has probability $1-e^{-\Omega(n)}$.
\end{proposition}
The proofs (\cref{app:seed}) combine Gaussian singular-value bounds \citep[Chapter~7]{vershynin2018high} with Weyl's inequality. The zero readout and the sign biases are essential: a random readout of scale $\varepsilon$ would start the loss of order $\varepsilon^2$ above its floor and destroy the exponential interval, and Gaussian biases would not fix the ReLU pattern.

\begin{theorem}[Conditional full-parameter process learning]\label{thm:process}
Let $\Theta(t)$ follow the gradient flow of $\Lproc$ on every parameter from \cref{eq:rand-seed}. Fix $\xi\in(0,1]$ and $\tau=\xi\log2/M$. Assume, on an event of probability at least $1-\zeta$ over the seed:
\begin{enumerate}[label=(P\arabic*),ref=(P\arabic*),leftmargin=*,itemsep=1pt]
\item\label{hyp:regular} \emph{Regularity.} An absolutely continuous solution exists on $[0,\infty)$ with $\frac{d}{dt}\Lproc=-\|\nabla_\Theta\Lproc\|^2$ almost everywhere.
\item\label{hyp:align} \emph{Persistent alignment.} There are $B_n,a_n>0$ and, whenever $\Lproc(\Theta(t))>\tau$, a direction $D(t)$ with
\begin{equation}\label{eq:align}
\|D(t)\|\le B_n,\qquad -\bigl\langle\nabla_\Theta\Lproc(\Theta(t)),D(t)\bigr\rangle\ge a_n\,\Lproc(\Theta(t)).
\end{equation}
\item\label{hyp:poly} \emph{Polynomial constants.} $B_n,\ a_n^{-1}\le\mathrm{poly}(n)$.
\end{enumerate}
Then $\Lproc(\Theta(t))\le\tau$ for all
\begin{equation}\label{eq:proc-time}
t\ \ge\ T_{\mathrm{proc}}:=\frac{M B_n^2}{a_n^2\,\xi\log2},
\end{equation}
and at every such time greedy decoding from a fresh prompt produces the entire gold continuation $s_1,\ldots,s_n,{:},s_n,\mathrm{EOS}$ with probability at least $1-\xi$. If $\xi^{-1}\le\mathrm{poly}(n)$, then $T_{\mathrm{proc}}=\mathrm{poly}(n)$.
\end{theorem}
\begin{proof}[Proof sketch]
By Cauchy--Schwarz, $a_nL\le\|\nabla L\|B_n$, so $\dot L=-\|\nabla L\|^2\le-(a_n/B_n)^2L^2$ and $L(t)<B_n^2/(a_n^2t)$ until $L\le\tau$; afterwards $L$ cannot increase. A wrong greedy token has probability at most $\frac12$, hence loss at least $\log2$. If free-running decoding errs, its first error is made on the gold prefix; a union bound over the $M$ positions and Markov's inequality give error probability at most $ML/\log2\le\xi$ (\cref{app:process}).
\end{proof}
$D(t)$ is only a witness: training follows the actual gradient of every tensor. No uniform transfer across histories is needed, because the guarantee is in probability over prompts.

\begin{proposition}[Readout separability implies alignment]\label{prop:separable}
Write $z_\Theta(x)=Vh_\Phi(x)$, with $h_\Phi(x)$ the final residual vector at context $x$ and $\Phi$ all parameters except $V$. Suppose that until $\Lproc\le\tau$, for almost every supervised context $x$ with gold token $y$,
\begin{enumerate}[label=(B\arabic*),ref=(B\arabic*),leftmargin=*,itemsep=1pt]
\item\label{hyp:sep} some $U(t)$ with $\|U(t)\|_F\le B_n$ has $(U_y-U_k)^\top h_{\Phi(t)}(x)\ge\gamma_n$ for every token $k\ne y$, and
\item\label{hyp:lossbd} the per-context loss is at most $H_n$.
\end{enumerate}
Then \ref{hyp:align} holds with $D=(U\text{ on }V,\ 0\text{ elsewhere})$ and $a_n=\gamma_n/(1+H_n)$, so
\begin{equation}\label{eq:sep-time}
T_{\mathrm{proc}}\le\frac{M B_n^2(1+H_n)^2}{\gamma_n^2\,\xi\log2}.
\end{equation}
\end{proposition}
\begin{proof}
$\nabla_V\ell=(p-e_y)h^\top$, so $-\langle\nabla_V\ell,U\rangle=\sum_{k\ne y}p_k(U_y-U_k)^\top h\ge\gamma_n(1-p_y)$, and $1-p_y=1-e^{-\ell}\ge\ell/(1+H_n)$. Average over contexts.
\end{proof}
\ref{hyp:sep} concerns the real evolving features; only the witness lives in readout coordinates. It is weak at the seed, where features are additive in the tokens up to score-order terms and cannot separate the transition label, so the readout witness certifies the phase after an interaction feature forms (\cref{app:process}).

\begin{corollary}[Matched separation from the random seed, conditional]\label{cor:separation}
Start both full flows from \cref{eq:rand-seed} with $\varepsilon=c_*n^{-1/2}$, widths of order $n$ and fixed $K,p$. If \ref{hyp:regular}--\ref{hyp:poly} hold for the process flow on the event of \cref{lem:seed} and $\xi^{-1}\le\mathrm{poly}(n)$, then for small $c_*$ and large $n$, $T_{\mathrm{proc}}<T_{\mathrm{out}}$: on $[T_{\mathrm{proc}},T_{\mathrm{out}}]$ the process model generates the whole continuation on all but a $\xi$ fraction of prompts while the outcome model satisfies \cref{eq:plateau}.
\end{corollary}
Both flows are Euclidean on all parameters, so no metric comparison enters. What remains is \ref{hyp:align} from the seed. \Cref{app:matrix} gives the exact unrestricted query--key flow $\dot M=-(W_Q^\top W_Q)N-N(W_K^\top W_K)$ for $M=W_Q^\top W_K$ and $N=\nabla_ML$, a conserved $W_QW_Q^\top-W_KW_K^\top$, and the score sensitivity through the ReLU MLP; they make explicit an interference term that the restricted family removes by construction. \Cref{app:cancel} states a conditional Gaussian cancellation lemma for that term, whose independence hypothesis does not yet hold for the evolving network, and \cref{app:obligations,app:diagnostics} list the open obligations and the measurements that test them.

\section{Experiments}\label{sec:experiments}
All numbers below summarize existing runs. We regenerate figures and tables from stored per-seed measurements; we do not present them as newly rerun training. Error bars describe seed variation, not uncertainty across independently drawn corpora. Protocols, evidence files and reproduction commands are in \cref{app:experiments}.

\paragraph{Controlled counting.}
Two-block, fixed-ReLU, small-initialization runs use $K=2$, $p=1/4$, $\varepsilon=0.002$, zero readout, exact full-batch float64 gradients over all words through $n=14$, and refined finite updates approximating flow. \Cref{fig:controlled} connects displacement times to $\lambda^{-n}$. It distinguishes exit from the theorem-radius ball from displacement to a larger diagnostic radius. The fast step exit fits the state-only prior; it is not a measurement of completed transition learning. These runs train unstructured seeds and test the barrier of \cref{thm:outcome}, not \cref{thm:step}.
\begin{figure}[t]
\centering\includegraphics[width=\linewidth]{figures/controlled.pdf}
\caption{\textbf{Outcome displacement times grow with horizon.} Stored runs, seeds $0,1,2$; medians and full seed ranges. Left: exit from $r_0\varepsilon$; right: displacement $0.1\varepsilon\|\tilde\theta\|$, beyond the theorem ball. Dashed guides are $\lambda^{-n}$ with $\lambda=1/2$, anchored at $n=4$. Step exit mainly learns the prior.}\label{fig:controlled}
\end{figure}
Actual transition learning is tested separately by finite-step routing runs. With query--key matrices, relative bias and readout trained from fan-in random weights, step training solves $16/17$ runs within $300$k updates (the remaining run reaches $99.95\%$); outcome solves $0/17$. At $n=6,8,12,16$, step solves $5/5,5/5,4/5,2/2$, with median solve updates $44$k,$61$k,$77$k,$237$k. These runs fix values, embeddings and MLPs but depart from the structured seed and supervise the first state too. Separately, the ten-seed AdamW comparison at $n=16$ reports $95.45\%$ mean step answer accuracy ($93.11\%$ complete-continuation accuracy) and $49.82\%$ outcome answer accuracy. They are empirical evidence outside the theorem, not time-bound verification.

\paragraph{Pretrained Boolean-8.}
Boolean-8 applies eight randomly wired NOT, CNOT, SWAP or Toffoli gates to four bits; there are $16$ answer states. Rank-$8$ LoRA on Qwen2.5-0.5B \citep{qwen2025qwen25,hu2022lora} reaches approximately $98.3\%$ held-out answer accuracy under step supervision and $11.1\%$ under outcome supervision. \Cref{tab:transfer} reports the checked seed summaries, the eightfold outcome budget, and answer-first control. These are synthetic sequential-computation results with a pretrained language model.
\input{figures/transfer_table}

\paragraph{BIG-Bench Hard object tracking.}
Training examples are generated from benchmark templates using fresh assignments and swaps. Testing uses the official $250$ held-out items at each of three, five and seven objects \citep{suzgun2022challenging}. Traces give every holder's item and its answer-option letter. Step supervision usually helps at three and seven objects. The five-object condition reverses the pattern: outcome achieves $80$--$94\%$, and one step seed reaches only $46\%$. \Cref{tab:transfer} retains all three sizes and all seeds. This reversal limits transfer of the mechanism: counting's contraction coefficient and the theorem's initialization do not characterize pretrained-model optimization on object tracking.

\paragraph{Count-label control.}
A separate ten-seed counting control matches one randomly selected supervised count label per prompt to the outcome count-label budget, with the same formatting labels and $30$k updates. Sparse process training attains mean answer accuracy $97.1,99.8,99.2\%$ for $n=6,8,12$ and solves $27/30$ runs at a $99\%$ threshold; outcome means are $68.5,64.3,51.2\%$, with $0/30$ solved. Matching labels does not match teacher-forced prefix information: sparse process inputs still contain gold running states. The shorter $8$k control was inconclusive and is retained in \cref{app:label-budget}.


\section{Related work}\label{sec:related}
\citet{kim2024parity} prove efficient teacher-forced parity learning with a Transformer and compare learning with and without chain of thought. Their parity setting, prescribed nonlinearity and attention parameterization differ from this paper's biased modular increments, local fixed-ReLU analysis and supplied-routing family. Our contribution is the matched finite-time plateau and persistent transition learning in the specified two-block softmax/ReLU parameterization; it does not subsume parity learning at $p=1/2$.

Makkuva and collaborators analyze Transformer loss landscapes and initialization-dependent dynamics on Markov data \citep{makkuva2025attention,makkuva2024local}. Their Markov next-token setting motivates an architecture-level optimization analysis. Here, the chain is the hidden modular state, and the comparison changes whether its intermediate states are displayed. We use Fourier conditioning rather than identify a separate attracting local minimum.

Gradient barriers based on weak target correlations \citep{shalevshwartz2017failures,abbe2025parities} explain why accessible features can reveal little about a high-order answer. The additional work here is controlling letter support, Taylor remainders and parameter scale in the Transformer, then converting a local loss floor into a trajectory bound. Automata constructions \citep{liu2023automata} show how shallow Transformers can represent sequential shortcuts; they do not imply convergence within the trained subspace. Scratchpad supervision \citep{nye2021scratchpads}, task decomposition \citep{wies2023subtask}, and locality of experience \citep{prystawski2023think} provide related perspectives on intermediate states. Expressivity and statistical advantages of chain of thought \citep{merrill2024expressive,wen2024sparse} concern complementary questions.


\section{Limitations and conclusion}\label{sec:conclusion}
The proved learning-time separation concerns population gradient flow, a sufficiently small structured initialization, fixed feature maps, aligned tied query--key scalars and supplied routing. The barrier is broader: it holds for gradient flow on every parameter from any small seed satisfying \cref{ass:init}, including a Gaussian seed with a weak perturbation, zero readout and sign biases. It relies on a fixed local ReLU pattern and lasts for a specified initial interval; it is not permanent failure, and it bounds the expected probability of the correct answer, not greedy accuracy.

\Cref{thm:process} removes the restrictions of the positive result and adds free-running generation, but it is conditional: it reduces full-parameter learning to persistent gradient alignment, and we do not prove that alignment from a random seed. Its guarantee holds with high probability over prompts, not on every prompt, and concerns predictions, not recovery of a particular routing or lookup mechanism. It is a statement about continuous-time population flow; step counts and sample sizes need separate arguments. The executors establish expressivity only. Empirical counting and pretrained-model runs change initialization, optimizer or architecture, and the five-object reversal prevents a universal transfer claim. Within the proved regime, however, the same trainable Transformer learns the recurring transition in polynomial time under displayed-state supervision while final-answer supervision remains near chance in expected correct-answer probability for an exponentially long initial interval.
